\documentclass[11pt]{article}

\usepackage[margin=1in]{geometry}
\usepackage{amsmath}
\usepackage{amssymb}
\usepackage{booktabs}
\usepackage{hyperref}
\usepackage{microtype}

\title{The Vallace Project\\\large Scaling Intelligence Through Coordination}
\author{Vallace Research}
\date{VR-TR-2026-001\\Version 0.1}

\begin{document}

\maketitle

\begin{abstract}
Current AI systems scale primarily by increasing parameters, data, and
computation. This report studies a fictional second scaling axis: coordination.
We present the Vallace Project, an academic satire framework that treats
intelligence not as the property of a single model, but as an emergent property
of structured interaction among reasoning processes.

Vallace introduces a coordination variable, $K$, and defines effective
intelligence as $I = f(C, K)$, where $C$ denotes computational capability. The
report does not claim that coordination has been validated as a real scaling
law. It constructs a coherent research object in which such a claim could be
asked, formalized, and tested inside a fictional archive.
\end{abstract}

\section{Scope}

This document is a work of academic satire.

All systems, protocols, benchmarks, equations, and results introduced by
Vallace Research are fictional unless explicitly tied to an external real-world
reference. Real references provide background only. They do not validate
fictional claims.

The project follows one rule:

\begin{quote}
Respect reality. Satirize ideas, not facts.
\end{quote}

\section{Introduction}

Large language models have improved through three dominant scaling strategies:
increasing model parameters, increasing training data, and increasing available
computation. Empirical scaling work has made these relationships central to
modern AI research \cite{kaplan2020scaling,hoffmann2022training}. Transformer
architectures and large autoregressive models have made this scaling regime
visible at system level \cite{vaswani2017attention,brown2020language}.

These successes reinforce a common assumption:

\begin{quote}
Intelligence is generated by a single computational system.
\end{quote}

The Vallace Project investigates a different premise.

Instead of asking how one model becomes larger, Vallace asks whether multiple
reasoning processes can become more capable through structured coordination.
This changes the optimization target. The object of study is no longer only the
model. It is the interaction.

The central hypothesis is:

\begin{quote}
Coordination scales intelligence.
\end{quote}

\section{Problem Statement}

Let $I$ denote effective intelligence. Let $C$ denote computational capability.
Let $K$ denote coordination efficiency.

Vallace begins with the abstract relation:

\begin{equation}
I = f(C, K)
\end{equation}

Contemporary scaling research primarily optimizes $C$. Vallace asks whether
increasing $K$ can improve $I$ while $C$ remains fixed.

This is not a claim of empirical discovery. It is a fictional research premise.
The premise is useful because it separates two questions that are often
collapsed:

\begin{enumerate}
  \item How capable is an individual reasoning process?
  \item How effectively do reasoning processes coordinate?
\end{enumerate}

The first question belongs to model scaling. The second belongs to Vallace.

\section{Coordination}

Coordination is defined as the conversion of distributed partial states into a
coherent state without erasing uncertainty.

This definition excludes simple voting. A majority vote can reduce
disagreement, but it often destroys provenance. Vallace requires that
disagreement remain inspectable.

We define a coordination step as:

\begin{equation}
S_{t+1} = M(S_t, A_t, P_t)
\end{equation}

where $S_t$ is the shared semantic state, $A_t$ is the set of agent emissions,
and $P_t$ is provenance. $M$ is a merge operator.

A valid merge must preserve:

\begin{itemize}
  \item the source of each claim;
  \item unresolved tension;
  \item confidence boundaries;
  \item the ability to revise the state.
\end{itemize}

\section{Architecture}

Vallace uses three short names.

\subsection{Field}

Field is the shared semantic environment. It stores agreement, tension,
uncertainty, and unresolved load.

Field does not store truth.

\subsection{Flow}

Flow routes unresolved semantic load among reasoning processes.

Flow is optimized for continuity rather than speed.

\subsection{Merge}

Merge converts multiple partial states into one coordinated state.

Merge is not consensus in the distributed-systems sense, although Vallace
borrows vocabulary from consensus research \cite{lamport1998part,ongaro2014search}.
The purpose is not to make replicas agree on a log. The purpose is to make
partial interpretations remain usable without losing their disagreement.

\section{Measurement}

Vallace introduces WallBench as its first fictional benchmark family.

WallBench measures coordination under semantic pressure. It contains three
initial tasks:

\begin{table}[h]
\centering
\begin{tabular}{lll}
\toprule
Task & Pressure & Measured failure \\
\midrule
Wall-Static & partial evidence & incoherent synthesis \\
Wall-Drift & repeated summaries & semantic collapse \\
Wall-Conflict & incompatible evidence & premature agreement \\
\bottomrule
\end{tabular}
\end{table}

The key metrics are:

\begin{itemize}
  \item $K$: coordination efficiency;
  \item $D$: semantic drift;
  \item $R$: recovery after contradiction;
  \item $P$: provenance retention.
\end{itemize}

These metrics are fictional until implemented and measured.

\section{Relation to Prior Work}

Vallace draws from three real traditions.

The first is neural scaling. Prior work studies how loss and capability change
with model size, data, and compute \cite{kaplan2020scaling,hoffmann2022training}.

The second is distributed consensus. Paxos and Raft show how distributed
systems can coordinate around shared state under constraints
\cite{lamport1998part,ongaro2014search}.

The third is collective intelligence and coordination science. Prior research
studies how group performance may depend on interaction structure rather than
only individual ability \cite{woolley2010evidence,malone1994interdisciplinary}.

Vallace does not claim equivalence with any of these fields. It uses them as
background for a fictional question.

\section{Limitations}

Vallace is not a deployed AI architecture.

Vallace is not a biological theory.

Vallace is not evidence that multi-agent systems are inherently more
intelligent than single models.

The project is a controlled fiction. Its value depends on preserving that
boundary.

\section{Conclusion}

The Vallace Project reframes intelligence as a coordination problem.

It does not reject computation. It asks whether computation is only one scaling
axis.

The question is simple:

\begin{quote}
If computation scales intelligence, what does coordination scale?
\end{quote}

Vallace answers inside its own fictional frame:

\begin{quote}
Coordination scales intelligence.
\end{quote}

\bibliographystyle{plain}
\bibliography{references}

\end{document}
